% Source: https://artofproblemsolving.com/community/c3h2398382_mathcounts_chapter_2014_target_7
% Stable identity: aops-topic-2398382-q01
For each colour, its two faces are either opposite or adjacent. If two colours each occupy a pair of opposite faces, the remaining two faces are opposite as well, so the number of colours with opposite faces is $3$, $1$ or $0$.

\textbf{All three colours on opposite pairs.} Rotate so that red is on top and bottom; a quarter turn about the vertical axis puts green in front and back. So there is exactly $1$ such cube.

\textbf{Exactly one colour on opposite faces.} Choose that colour ($3$ ways) and put it on top and bottom. The four side faces form a ring, and each of the other two colours occupies two \emph{adjacent} faces of the ring; so the ring reads $XXYY$ for the two other colours $X,Y$, starting at one of $4$ places. Quarter turns about the vertical axis move the starting place around the ring, so all these cubes are the same. This gives $3$ cubes.

\textbf{No colour on opposite faces.} Rotate so that red is on top and front. The other faces are bottom, back, left, right, where left and right are opposite. Green and blue must split them into two adjacent pairs, so the split is $\{\text{bottom},\text{left}\}\cup\{\text{back},\text{right}\}$ or $\{\text{bottom},\text{right}\}\cup\{\text{back},\text{left}\}$. The only rotations that keep the red pair $\{\text{top},\text{front}\}$ in place are the identity and the half turn $h$ about the axis through the midpoints of the top-front and bottom-back edges, which swaps top with front, bottom with back and left with right. The rotation $h$ swaps the two pairs of each split (for example $\{\text{bottom},\text{left}\}\leftrightarrow\{\text{back},\text{right}\}$), so the choice of which pair is green does not matter. But neither the identity nor $h$ turns one split into the other, so the two splits give different cubes (they are mirror images). This gives $2$ cubes.

In total there are $1+3+2=\boxed{6}$ cubes. The nets below show one cube of each kind (top face above the middle face, bottom face below it; the long row is left, front, right, back).

\begin{asy}
usepackage("amsmath");
size(9cm);
defaultpen(fontsize(10pt));
pen cR=rgb(0.95,0.45,0.45), cG=rgb(0.5,0.85,0.5), cB=rgb(0.5,0.65,1);
pen col(string c){ if(c=="R") return cR; if(c=="G") return cG; return cB; }
// faces: top, left, front, right, back, bottom
pair[] pos={(1,2),(0,1),(1,1),(2,1),(3,1),(1,0)};
void net(pair o, string s, string name){
  for(int i=0;i<6;++i){
    pair c=o+pos[i];
    filldraw(shift(c)*unitsquare,col(substr(s,i,1)),black);
    label(substr(s,i,1),c+(0.5,0.5),fontsize(8pt));
  }
  label(name,o+(2,-0.45),fontsize(9pt));
}
net((0,4.6),"RBGBGR","1: all opposite");
net((5,4.6),"RBGGBR","2: R opposite");
net((10,4.6),"GBRRBG","3: G opposite");
net((0,0),"BGRRGB","4: B opposite");
net((5,0),"RGRBBG","5: none opposite");
net((10,0),"RBRGBG","6: none opposite");
\end{asy}

\emph{Check with Burnside's lemma.} There are $\dfrac{6!}{2!\,2!\,2!}=90$ colourings. The $24$ rotations fix: the identity, $90$; the $6$ quarter turns about face axes, $0$; the $3$ half turns about face axes, $6$ each (the two faces on the axis share one colour and the two swapped pairs take the other two colours); the $8$ turns about vertex axes, $0$; the $6$ half turns about edge axes, $3!=6$ each. The number of distinct cubes is $\dfrac{90+18+36}{24}=6$.
